\documentclass[11pt,a4paper]{article}

\usepackage[utf8]{inputenc}
\usepackage[T1]{fontenc}
\usepackage[brazilian]{babel}
\usepackage[margin=2cm]{geometry}
\usepackage{enumitem}
\usepackage{hyperref}
\usepackage{titlesec}
\usepackage{parskip}

\hypersetup{
  colorlinks=true,
  linkcolor=black,
  urlcolor=black,
  pdftitle={Kamille Hartmann Maccagnan - Currículo},
  pdfauthor={Kamille Hartmann Maccagnan}
}

\titleformat{\section}{\large\bfseries}{}{0em}{}[\titlerule]
\titlespacing*{\section}{0pt}{12pt}{6pt}

\pagestyle{empty}
\setlength{\parindent}{0pt}
\setlength{\parskip}{6pt}

\newcommand{\cvheader}[4]{%
  \begin{center}
    {\LARGE\bfseries #1}\\[4pt]
    {\large #2}\\[6pt]
    #3\\[2pt]
    #4
  \end{center}
  \vspace{4pt}
}

\newcommand{\experience}[3]{%
  \textbf{#1} \hfill \textit{#2}\\[4pt]
  #3
  \vspace{8pt}
}

\begin{document}

\cvheader{Kamille Hartmann Maccagnan}%
{Consultora de Serviços de TI | Onboarding | Consultoria Funcional | Gestão de Projetos}%
{Campo Comprido, Curitiba-PR \textbullet{} (41) 9 9928-4424 \textbullet{} \href{mailto:kamillehartmannm@gmail.com}{kamillehartmannm@gmail.com}}%
{\href{https://linkedin.com/in/kamille-hartmann-maccagnan/}{linkedin.com/in/kamille-hartmann-maccagnan}}

\section*{Resumo Profissional}

Profissional com experiência em onboarding, consultoria funcional e gestão de demandas em ambiente de tecnologia e operações. Atuação na condução de treinamentos, implantação de processos e relacionamento com usuários como ponto focal de T.I., com gestão de interações por WhatsApp corporativo em conversas individuais e grupos. Vivência em homologação de sistemas, elaboração de documentação funcional, POPs e materiais de apoio, promovendo adoção, autonomia e qualidade na operação. Experiência em PMO com acompanhamento de cronogramas, marcos de entrega, análise de riscos, identificação de pendências e alinhamento com stakeholders em ambiente multinacional. Perfil analítico, organizado e comunicativo, com monitoramento de indicadores em Power BI, priorização de múltiplas demandas e inglês avançado.

\section*{Experiência Profissional}

\experience{Anjuss -- Analista de Suporte de TI}{07/2025 -- Atual}{%
Atuação na condução da jornada de entrada e capacitação de novos funcionários, com liderança de treinamentos, onboarding e disseminação de conhecimento entre lojas, operações e T.I. Responsável pela gestão de interações via WhatsApp corporativo, conduzindo o alinhamento entre operação e tecnologia de forma consultiva: entendo a necessidade do usuário, oriento quando o tema é operacional e escalo para o time do sistema ou para adquirentes REDE/Cielo quando necessário, acompanhando o retorno até a resolução e promovendo autonomia no uso da solução. Participação em testes em ambiente de homologação do sistema de vendas, com validação de fluxos antes da liberação para as lojas, garantindo critérios de qualidade e ativação da operação. Elaboração de documentação funcional estruturada, incluindo POPs e materiais de apoio, com registro e rastreabilidade das atividades e evolução das demandas. Desenvolvimento de dashboards em Power BI para monitoramento de indicadores, identificação proativa de riscos operacionais e priorização de ações.}

\experience{HORSE do Brasil -- Estagiária de TI (IS/IT LATAM)}{07/2024 -- 07/2025}{%
Atuação em PMO e gestão de projetos em ambiente multinacional, com acompanhamento de cronogramas, marcos de entrega, indicadores de desempenho e visibilidade de status de iniciativas. Elaboração de relatórios executivos e materiais de acompanhamento para apresentação de resultados e alinhamento com stakeholders de negócio, fornecedores e times técnicos, com participação em reuniões conduzidas em até três idiomas simultaneamente. Atuação na documentação de processos, análise de riscos, identificação de pendências, investigação de desvios operacionais e facilitação de comunicação entre múltiplas áreas para garantir prazos, entregáveis e conformidade com combinados. Participação em iniciativas de melhoria contínua e evolução de práticas de gestão com foco em qualidade de entregas.}

\experience{Restaurante Familiar -- Analista de Dados e Suporte Técnico}{01/2023 -- 07/2024}{%
Atuação no suporte à gestão do negócio com análise de performance operacional e resultados financeiros. Estruturação de base de dados em Excel e dashboards em Power BI para acompanhamento de indicadores, com padronização de informações e documentação de procedimentos. Responsável por diagnosticar inconsistências, propor melhorias operacionais e traduzir necessidades do negócio em rotinas práticas de acompanhamento e validação, com atuação no alinhamento entre áreas operacionais e administrativas.}

\section*{Projetos e Conquistas}

Estruturação de KPIs e dashboards para monitoramento de indicadores e suporte à tomada de decisão. Participação em iniciativas de melhoria contínua e otimização de processos. 1º lugar no Transformation Day Renault 2023 com solução orientada a dados e melhoria de processos. Participação no Hackathon Bosch 2023 com foco em resolução de problemas de negócio com uso de dados.

\section*{Habilidades Técnicas}

\textbf{Onboarding e Consultoria Funcional:} Condução de treinamentos e onboarding, atendimento consultivo, documentação funcional, POPs, adoção de sistemas, gestão de WhatsApp corporativo, repasse de conhecimento, tradução de necessidades de negócio em orientações práticas.

\textbf{Gestão de Projetos e Demandas:} PMO, acompanhamento de cronogramas e marcos, análise de riscos, priorização de ações, remoção de impedimentos, rastreabilidade de atividades, Scrum, Kanban, Boas Práticas PMOBOK.

\textbf{Análise e Indicadores:} Diagnóstico de problemas, KPIs, dashboards, Power BI (DAX, Power Query), Excel, SQL, monitoramento de performance, melhoria contínua.

\textbf{Ferramentas:} Monday, Trello, ServiceNow, Spotfire, Pacote Office.

\textbf{Comunicação e Relacionamento:} Interface com stakeholders, comunicação clara, autonomia, organização, inglês avançado.

\textbf{Idiomas:} Inglês avançado \textbullet{} Espanhol básico \textbullet{} Coreano básico.

\section*{Formação Acadêmica}

\textbf{PUCPR} -- Graduação em Gestão da Tecnologia da Informação \hfill Previsão de conclusão: 06/2026\\
Curitiba, PR

\textbf{Escola Conquer} -- Pós-graduação em Liderança e Gestão em Tecnologia \hfill Conclusão: 2024\\
Curitiba, PR

\textbf{Universidade Tuiuti do Paraná} -- Graduação em Medicina Veterinária \hfill Conclusão: 06/2019\\
Curitiba, PR

\textbf{Qualittas} -- Pós-graduação em Clínica Médica e Cirúrgica de Animais Exóticos e Pets Não Convencionais \hfill Conclusão: 12/2022\\
Curitiba, PR

\end{document}
